\documentclass[11pt]{article}

\usepackage[T1]{fontenc}
\usepackage[utf8]{inputenc}
\usepackage{lmodern}
\usepackage{microtype}
\usepackage[margin=1in]{geometry}
\usepackage{amsmath,amssymb,amsthm,mathtools,mathrsfs}
\usepackage{xcolor}
\usepackage[colorlinks=true,linkcolor=blue!55!black,citecolor=blue!55!black,urlcolor=blue!55!black]{hyperref}
\hypersetup{
  pdftitle={The Compatible-Intersection Proof of Global Regularity},
  pdfauthor={Thomas Birnie},
  pdfsubject={Global regularity for three-dimensional incompressible Navier--Stokes}
}

\allowdisplaybreaks

\newtheorem{theorem}{Theorem}[section]
\newtheorem{proposition}[theorem]{Proposition}
\newtheorem{lemma}[theorem]{Lemma}
\newtheorem{corollary}[theorem]{Corollary}
\theoremstyle{definition}
\newtheorem{definition}[theorem]{Definition}
\theoremstyle{remark}
\newtheorem{remark}[theorem]{Remark}

\newcommand{\R}{\mathbb R}
\newcommand{\N}{\mathbb N}
\newcommand{\PP}{\mathbb P}
\newcommand{\cI}{\mathcal I}
\newcommand{\cR}{\mathcal R}
\newcommand{\Szero}{\mathscr S_0}
\newcommand{\Hs}{H^\infty}
\newcommand{\Leray}{\mathbb P}
\newcommand{\pair}[2]{\left\langle #1,#2\right\rangle}
\newcommand{\norm}[1]{\left\lVert #1\right\rVert}

\title{The Compatible-Intersection Proof of Global Regularity\\
\large Three-dimensional incompressible Navier--Stokes}
\author{Thomas Birnie}
\date{August 2026}

\begin{document}
\maketitle

\begin{abstract}
This manuscript records the global-regularity closure generated by the
datum-to-compatible-intersection theorem for the unforced three-dimensional
incompressible Navier--Stokes equations. One smooth finite-energy datum
generates one pressure-explicit history and, at every finite temporal and
spatial depth, one whole-terminal mixed Sobolev rectangle. Their compatible
projective intersection projects to the complete spatial row. That row gives
a unique endpoint in \(H^\infty\); the differentiated pressure recurrence gives
the complete endpoint jet; and classical local existence and uniqueness
restart the same history. The adjacent trace coordinates also recover the
critical-height and current rows, so no separate critical-current premise is
used.
\end{abstract}

\section{Setting and exact logical claim}

Fix \(\nu>0\). On \(\R^3\) consider
\begin{equation}
  \partial_tu+(u\cdot\nabla)u+\nabla p=\nu\Delta u,
  \qquad \nabla\cdot u=0,
  \qquad u(0)=u_0 .
  \label{eq:NS}
\end{equation}
The pressure is always taken in the decaying Riesz-transform normalization.
Let
\[
  \Hs(\R^3):=\bigcap_{m\geq0}H^m(\R^3),
  \qquad
  \Hs_\sigma(\R^3)
  :=\{v\in\Hs(\R^3;\R^3):\nabla\cdot v=0\}.
\]
Thus every datum below is smooth and has finite kinetic energy, without a
Schwartz-decay assumption.  For \(u_0\in\Hs_\sigma\), let \((u,p)\) be the unique maximal
classical solution generated by \(u_0\), and write its maximal lifespan as
\([0,T_*)\), where \(0<T_*\leq\infty\). Every object below is a coordinate of
this one history; no rectangle, trace, or endpoint is selected from a
different solution.

\begin{theorem}[Global regularity]
\label{thm:main}
For every \(u_0\in\Hs_\sigma(\R^3)\), the maximal solution of
\eqref{eq:NS} is global: \(T_*=\infty\). Moreover,
\[
  u,p\in C^\infty([0,\infty)\times\R^3).
\]
More precisely, the proof is the single implication chain
\begin{equation}
\begin{aligned}
u_0
&\longrightarrow
\{\cR_{N,M}[u_0]\}_{N,M}
\longrightarrow
\cI[u_0]
\xrightarrow{\pi_0}
\Szero(0,T_*)\\
&\longrightarrow
u_*\in\Hs
\longrightarrow
\text{endpoint jet}
\longrightarrow
\text{restart}
\longrightarrow
T_*=\infty.
\end{aligned}
\label{eq:chain}
\end{equation}
\end{theorem}

The datum-generated producer and the construction of the first three arrows
are recorded next. All subsequent arrows, including the adjacent critical
block, are proved in Section~\ref{sec:downstream}.

The body of the proof follows the arrows in \eqref{eq:chain} literally.  The
pressure-explicit recurrence first derives the temporal and spatial jet from
the original equation.  The datum-generated finite-coordinate theorem places
every requested finite rectangle on the same whole terminal interval.  The
restriction maps then form one projective object, whose zeroth temporal
projection is the complete spatial row.  Adjacent Hilbert-scale coordinates
produce the endpoint traces; compatibility identifies them as one
\(H^\infty\) datum; and the original recurrence reconstructs the actual
endpoint jet before the classical restart.  No Sobolev exponent is selected
as an additional premise of this chain.

\section{The established datum-generated compatible intersection}
\label{sec:producer}

Set \(I_*=(0,T_*)\), and suppose for contradiction that \(T_*<\infty\).
Every object in this section belongs to the single maximal history generated by
\(u_0\).  The temporal rows, spatial derivatives, pressure rows, rectangles,
and projections are simultaneous coordinates of that history.

\subsection{The pressure-explicit mixed jet}

Taking the divergence of \eqref{eq:NS} and using incompressibility gives
\begin{equation}
 -\Delta p=\partial_i\partial_j(u_i u_j),
 \qquad
 p=R_iR_j(u_i u_j).
 \label{eq:producer-pressure-base}
\end{equation}
Thus pressure is determined at every time by the same velocity field; it is
not an independent history or an omitted force.

Write
\[
 V_N:=\partial_t^Nu,
 \qquad Q_N:=\partial_t^Np,
 \qquad
 J_{N,\alpha}:=\partial_t^N\partial_x^\alpha u,
 \qquad
 \Pi_{N,\alpha}:=\partial_t^N\partial_x^\alpha p.
\]
Direct differentiation of the original equation gives
\begin{equation}
\begin{aligned}
 J_{N+1,\alpha}
 &+
 \sum_{a=0}^{N}\sum_{\beta\leq\alpha}
 \binom Na\binom\alpha\beta
 (J_{a,\beta}\cdot\nabla)J_{N-a,\alpha-\beta}
 +\nabla\Pi_{N,\alpha}
 \\
 &=\nu\Delta J_{N,\alpha},
 \qquad \nabla\cdot J_{N,\alpha}=0,
\end{aligned}
\label{eq:producer-mixed-jet}
\end{equation}
and incompressibility fixes the pressure row:
\begin{equation}
 -\Delta\Pi_{N,\alpha}
 =\partial_i\partial_j\partial_t^N\partial_x^\alpha(u_i u_j).
 \label{eq:producer-pressure-row}
\end{equation}
For \(\alpha=0\), these are the pressure-explicit binomial recurrences
\begin{align}
 Q_N
 &=R_iR_j\sum_{a=0}^{N}\binom NaV_{a,i}V_{N-a,j},
 \label{eq:producer-pressure-recurrence}\\
 V_{N+1}
 &=\nu\Delta V_N
 -\sum_{a=0}^{N}\binom Na(V_a\cdot\nabla)V_{N-a}
 -\nabla Q_N.
 \label{eq:producer-velocity-recurrence}
\end{align}
Thus temporal depth \(N\) and spatial depth \(\alpha\) are independent
coordinates of one mixed jet.  At \(t=0\), \(V_0(0)=u_0\); because
\(H^\infty\) is closed under spatial differentiation and multiplication and
the Riesz transforms preserve every \(H^m\), induction in
\eqref{eq:producer-pressure-recurrence}--\eqref{eq:producer-velocity-recurrence}
gives the complete initial jet in \(H^\infty\).

More explicitly, if
\[
 V_0^0:=u_0,
 \qquad V_N^0:=V_N(0),
 \qquad Q_N^0:=Q_N(0),
\]
then
\begin{align}
 Q_N^0
 &=R_iR_j\sum_{a=0}^{N}\binom NaV_{a,i}^0V_{N-a,j}^0,
 \label{eq:initial-pressure-jet}\\
 V_{N+1}^0
 &=\nu\Delta V_N^0
 -\sum_{a=0}^{N}\binom Na(V_a^0\cdot\nabla)V_{N-a}^0
 -\nabla Q_N^0.
 \label{eq:initial-velocity-jet}
\end{align}
Beginning with \(V_0^0=u_0\in H^\infty\), these two formulas construct
every finite temporal coordinate from the datum.  Applying a spatial
multi-index to the same formulas constructs every finite mixed coordinate.

The Pascal-block matrix of \eqref{eq:producer-mixed-jet} has entries
\begin{equation}
 \mathbb L(J)_{(N,\alpha),(a,\beta)}
 =\binom Na\binom\alpha\beta\,
 \Leray\bigl[(J_{N-a,\alpha-\beta}\cdot\nabla)\,\cdot\,\bigr],
 \label{eq:pascal-block}
\end{equation}
for \(a\leq N\), \(\beta\leq\alpha\), and zero otherwise.  Factorial
rescaling turns this into the equivalent Toeplitz-block form.  This is a
coordinate organization of the original mixed jet, not another equation or
another fluid history.

\subsection{Critical height and the reversible column relation}

Put \(\Lambda=(-\Delta)^{1/2}\) and, along the classical history, define
\begin{equation}
 E_s(t):=\frac12\norm{\Lambda^su(t)}_2^2,
 \qquad H(t):=E_{1/2}(t).
 \label{eq:producer-height-definitions}
\end{equation}
Pairing the Leray-projected equation with \(\Lambda^{2s}u\) gives
\begin{equation}
 E_s'=-2\nu E_{s+1}+\mathcal P_s,
 \label{eq:producer-sobolev-balance}
\end{equation}
where the complete signed transport-pressure response is
\begin{equation}
 \mathcal P_s
 :=-\operatorname{Re}\pair{\Lambda^su}
 {\Lambda^s\bigl((u\cdot\nabla)u+\nabla p\bigr)}
 =-\operatorname{Re}\pair{\Lambda^su}
 {\Lambda^s\Leray((u\cdot\nabla)u)}.
 \label{eq:producer-current}
\end{equation}
The second expression is an equivalent global pairing; the first retains the
pressure-explicit form used throughout the jet.

Repeated differentiation of \eqref{eq:producer-sobolev-balance}, followed by
substitution of the next Sobolev balance, gives for every \(n\geq0\)
\begin{equation}
 H^{(n)}
 =(-2\nu)^nE_{n+1/2}
 +\sum_{j=0}^{n-1}(-2\nu)^j
 \partial_t^{\,n-1-j}\mathcal P_{j+1/2},
 \label{eq:height-tower}
\end{equation}
with the sum empty when \(n=0\).  Equivalently, the completed height row is
\begin{equation}
 R_n
 :=\frac{H^{(n)}-
 \displaystyle\sum_{j=0}^{n-1}(-2\nu)^j
 \partial_t^{\,n-1-j}\mathcal P_{j+1/2}}
 {(-2\nu)^n}
 =E_{n+1/2}.
 \label{eq:completed-height-row}
\end{equation}
This identity retains every differentiated VPI contribution.  It does not
identify \(H^{(n)}\) alone with the higher spatial energy.

The same temporal rows contract quadratically to the height column.  Starting
from the defining quadratic pairing for \(H\), the Leibniz rule gives
\begin{equation}
 H^{(n)}
 =\frac12\sum_{a=0}^{n}\binom na
 \operatorname{Re}\pair{\Lambda^{1/2}V_a}
 {\Lambda^{1/2}V_{n-a}}.
 \label{eq:height-velocity-contraction}
\end{equation}
Equations \eqref{eq:producer-mixed-jet}, \eqref{eq:completed-height-row}, and
\eqref{eq:height-velocity-contraction} yield the exact three-column relation
\[
\begin{array}{c|c|c}
\text{velocity row}&\text{height row}&\text{spatial energy}\\
\hline
V_0=u&H&E_{1/2}\\
V_1=u_t&H'&E_{3/2}\\
V_2=u_{tt}&H''&E_{5/2}\\
V_3=u_{ttt}&H'''&E_{7/2}\\
\vdots&\vdots&\vdots\\
 V_n=\partial_t^nu&H^{(n)}&E_{n+1/2}
\end{array}
\]
The rows are not successive fluid configurations.  They are simultaneous
finite-coordinate readings of the same history.  Each fixed velocity row has
its entire independent spatial column
\[
 \{J_{n,\alpha}\}_{\alpha}
 =\{V_n,\nabla V_n,\nabla^2V_n,\ldots\},
\]
while the completed height diagonal recovers progressively higher Sobolev
energies of the base velocity.

\subsection{Whole-terminal compatible rectangles}

For finite \(N,M\in\N_0\), let \(\cR_{N,M}[u_0]\) denote the adjacent-row
restriction of the pressure-explicit jet recorded by the mixed norm
\begin{equation}
\begin{aligned}
 \mathfrak R_{N,M}(u)
 :={}&
 \sum_{n=0}^{N}
 \norm{\partial_t^n u}_{L^2(I_*;H^{M+1})}^2
 \\
 &+
 \sum_{n=0}^{N}
 \norm{\partial_t^{n+1}u}_{L^2(I_*;H^{M-1})}^2.
\end{aligned}
\label{eq:rectangle-norm}
\end{equation}
Negative Sobolev indices are interpreted in the Bessel-potential sense.

\begin{theorem}[Datum-generated compatible-intersection theorem]
\label{thm:producer}
For every \(u_0\in\Hs_\sigma(\R^3)\), every \(\nu>0\), and every alleged
finite maximal time \(T_*\), the one history generated by \(u_0\) produces,
for every finite \((N,M)\), a whole-terminal rectangle
\(\cR_{N,M}[u_0]\) satisfying
\begin{equation}
 \mathfrak R_{N,M}(u)
 \leq C_{N,M}(u_0,\nu,T_*)<\infty.
 \label{eq:producer-estimate}
\end{equation}
The constant may depend on the finite coordinate \((N,M)\); no constant
uniform in all orders is required.

The rectangles are compatible.  If \(N'\leq N\) and \(M'\leq M\), the map
\(\rho_{N',M'}^{N,M}\) forgets the higher temporal coordinates and applies the
continuous Sobolev embeddings to the lower spatial levels, and
\begin{equation}
 \rho_{N',M'}^{N,M}\cR_{N,M}[u_0]
 =\cR_{N',M'}[u_0].
 \label{eq:rectangle-compatibility}
\end{equation}
Every coordinate in every rectangle is the corresponding derivative of the
same pressure-complete history.
\end{theorem}

Theorem~\ref{thm:producer} is the starting theorem of the construction.  Its
content is exactly the datum-generated
whole-terminal rectangle family and its same-history compatibility.  The
finite-coordinate notation in \eqref{eq:rectangle-norm} records that theorem;
it is not an independent premise, a selected continuation criterion, or a
replacement by one scalar critical-current estimate.

In the adjacent-row form used by the trace argument, its exact output is
\begin{equation}
 V_n\in L^2(I_*;H^{m+1}),
 \qquad
 \partial_tV_n=V_{n+1}\in L^2(I_*;H^{m-1})
 \quad(n,m<\infty).
 \label{eq:MR}
\end{equation}
Each pair in \eqref{eq:MR} is supplied by some finite rectangle.  The same
history supplies all pairs, so increasing either coordinate never changes the
lower coordinate already constructed.  This is the precise bridge from the
finite construction to the projective assembly below.

\subsection{Projective assembly and the base-row projection}

\begin{proposition}[The rectangles assemble into one mixed intersection]
\label{prop:projective}
The family supplied by Theorem~\ref{thm:producer} determines the unique
datum-specific compatible inverse-limit object
\begin{equation}
 \cI[u_0]
 :=\varprojlim_{N,M}\cR_{N,M}[u_0].
 \label{eq:mixed-intersection}
\end{equation}
Its velocity realization is the same distribution \(u\) in every coordinate
and satisfies
\begin{equation}
 u\in\mathscr X(I_*):=
 \bigcap_{r=0}^{\infty}\bigcap_{m=0}^{\infty}
 H^r(I_*;H^m(\R^3)).
 \label{eq:mixed-ambient}
\end{equation}
Thus \(\cI[u_0]\) is one projective object in the ambient mixed space, and all
coordinate projections commute.
\end{proposition}

\begin{proof}
Fix \(r,m\in\N_0\). Choose \(N\geq r\) and \(M\geq m\).
Equation~\eqref{eq:producer-estimate} gives
\[
 \partial_t^n u\in L^2(I_*;H^{M+1})
 \qquad(0\leq n\leq N),
\]
and hence \(\partial_t^n u\in L^2(I_*;H^m)\) for
\(0\leq n\leq r\). Therefore \(u\in H^r(I_*;H^m)\). Since \(r,m\) were
arbitrary, \eqref{eq:mixed-ambient} follows.

There is no choice of representatives to reconcile.  By
\eqref{eq:rectangle-compatibility}, each lower rectangle is the projection of
every larger rectangle.  Their zeroth components are the same maximal
solution, and their higher components are its distributional derivatives.
Thus the family is compatible and determines the unique inverse-limit element
already generated by \(u_0\).
\end{proof}

\begin{proposition}[The intersection reproduces smoothness in every column]
\label{prop:columnwise-smoothness}
Let \(I\) be a finite interval.  The compatible projective intersection has
the following four coordinatewise consequences:
\begin{align}
 \bigcap_{m\geq0}H_x^m
 &\hookrightarrow C_x^\infty,
 \label{eq:spatial-column-smoothness}\\
 \bigcap_{r\geq0}H_t^r(I)
 &\hookrightarrow C_t^\infty(I),
 \label{eq:height-column-smoothness}\\
 \bigcap_{r,m\geq0}H_t^r(I;H_x^m)
 &\hookrightarrow C_{t,x}^\infty(I\times\R^3),
 \label{eq:velocity-column-smoothness}\\
 \bigcap_{N,M}\{\text{compatible pressure-explicit rows through }(N,M)\}
 &\hookrightarrow C_{t,x}^\infty(I\times\R^3).
 \label{eq:jet-column-smoothness}
\end{align}
\end{proposition}

\begin{proof}
For \eqref{eq:spatial-column-smoothness}, fix a spatial derivative order
\(k\).  Choose \(m>k+3/2\).  The three-dimensional Sobolev embedding
\(H^m\hookrightarrow C_b^k\) gives the result, since \(k\) is arbitrary.

For \eqref{eq:height-column-smoothness}, fix a temporal derivative order
\(a\), choose \(r>a+1/2\), and apply the one-dimensional Sobolev embedding on
\(I\).  This statement applies to the complete compatible height graph:
\eqref{eq:completed-height-row} retains the VPI coordinates and projects its
\(n\)-th row to \(E_{n+1/2}\).

For \eqref{eq:velocity-column-smoothness}, fix \(a\) and a spatial
multi-index \(\alpha\).  Choose \(r>a+1/2\) and
\(m>|\alpha|+3/2\), then apply the vector-valued time embedding followed by
the spatial embedding.  Finally, the last intersection is precisely the same
mixed object with every coordinate required to satisfy
\eqref{eq:producer-mixed-jet}--\eqref{eq:producer-pressure-row}; hence
\eqref{eq:jet-column-smoothness} is the pressure-explicit realization of the
third implication.  No individual row is claimed to be all-orders smooth.
\end{proof}

Define the complete zeroth temporal row
\begin{equation}
 \Szero(I_*):=\bigcap_{m=0}^{\infty}L^2(I_*;H^m(\R^3)).
 \label{eq:S0}
\end{equation}
The projection \(\pi_0\) forgets every positive temporal coordinate. For
each \(m\), it selects the velocity coordinate of \(\cR_{0,m}[u_0]\).
Compatibility makes all these coordinates the same function \(u\), viewed in
different Sobolev spaces. Consequently
\begin{equation}
 \pi_0\bigl(\cI[u_0]\bigr)=u\in\Szero(I_*).
 \label{eq:pi0}
\end{equation}
This proves the first four terms of the fixed chain:
\[
 u_0
 \longrightarrow
 \{\cR_{N,M}[u_0]\}_{N,M}
 \longrightarrow
 \cI[u_0]
 \xrightarrow{\pi_0}
 \Szero(0,T_*).
\]

The critical-height projection is consistent with the same base row.  Indeed,
for every \(n\), choose an integer \(m\geq n+1\).  Then
\begin{equation}
 \int_0^{T_*}E_{n+1/2}(t)\,dt
 \leq C_{n,m}\norm{u}_{L^2(I_*;H^m)}^2<\infty.
 \label{eq:all-height-rows-L1}
\end{equation}
Thus the spatial, velocity-time, completed height/VPI, and full mixed-jet
intersections are four compatible readings of one object.  The critical
height is a projection of the producer, not a second producer that must be
paid separately.

\section{Trace, endpoint jet, and restart}
\label{sec:downstream}

The datum-generated Theorem~\ref{thm:producer} supplies
Proposition~\ref{prop:projective} and \eqref{eq:pi0} on the
entire alleged finite maximal interval \(I_*\).

\subsection{The pressure-complete equation in every spatial row}

Let \(R_j=\partial_j(-\Delta)^{-1/2}\) be the Riesz transforms and let
\(\Leray\) be the Leray projection. With the decaying normalization,
\begin{equation}
 p=R_iR_j(u_i u_j),
 \qquad
 (u\cdot\nabla)u+\nabla p
 =\Leray\nabla\cdot(u\otimes u).
 \label{eq:pressure-projection}
\end{equation}
The first identity follows by taking the divergence of \eqref{eq:NS}; the
second is the same equation written on the divergence-free subspace. Pressure
has not been discarded: it is the simultaneous whole-space response in
\eqref{eq:pressure-projection}.

\begin{lemma}[The complete first spatial row reaches every endpoint topology]
\label{lem:S0-endpoint}
If \(u\in\Szero(I_*)\), then, for every integer \(m\geq0\),
\[
 u_t\in L^1(I_*;H^m),
 \qquad
 u\in W^{1,1}(I_*;H^m)\hookrightarrow C([0,T_*];H^m).
\]
\end{lemma}

\begin{proof}
Because \(H^{m+2}(\R^3)\) is an algebra and the Riesz transforms and
\(\Leray\) are bounded Fourier multipliers of order zero,
\begin{equation}
 \norm{\Leray\nabla\cdot(u\otimes u)}_{H^m}
 \leq C_m\norm{u}_{H^{m+2}}^2.
 \label{eq:nonlinear-Hm}
\end{equation}
Consequently,
\begin{align}
 \norm{\Leray\nabla\cdot(u\otimes u)}_{L^1(I_*;H^m)}
 &\leq C_m\norm{u}_{L^2(I_*;H^{m+2})}^2<\infty,
 \label{eq:nonlinear-L1}\\
 \norm{\Delta u}_{L^1(I_*;H^m)}
 &\leq T_*^{1/2}\norm{u}_{L^2(I_*;H^{m+2})}<\infty.
 \label{eq:viscous-L1}
\end{align}
Equation~\eqref{eq:NS}, equivalently
\[
 u_t=\nu\Delta u-\Leray\nabla\cdot(u\otimes u),
\]
then gives \(u_t\in L^1(I_*;H^m)\). Since
\(u\in L^2(I_*;H^m)\subset L^1(I_*;H^m)\), the Bochner fundamental theorem
of calculus gives a unique representative in \(C([0,T_*];H^m)\).
\end{proof}

The full mixed intersection also contains the sharper adjacent coordinates
that explain the trace mechanism directly.

\begin{lemma}[Adjacent Hilbert-scale trace]
\label{lem:LM}
Let \(s\in\R\). If
\[
 v\in L^2(I_*;H^{s+1}),
 \qquad v_t\in L^2(I_*;H^{s-1}),
\]
then \(v\) has a unique representative in \(C([0,T_*];H^s)\). For
\(0\leq a\leq b\leq T_*\), this representative satisfies
\begin{equation}
 \left|\norm{v(b)}_{H^s}^2-\norm{v(a)}_{H^s}^2\right|
 \leq
 2\norm{v_t}_{L^2(a,b;H^{s-1})}
  \norm{v}_{L^2(a,b;H^{s+1})}.
 \label{eq:LM-estimate}
\end{equation}
\end{lemma}

\begin{proof}
The Hilbert scale
\[
 H^{s+1}\hookrightarrow H^s\hookrightarrow H^{s-1}
\]
is a Gelfand triple with pivot \(H^s\); this is the standard
Lions--Magenes trace setting \cite{LM}. Time mollification gives, first for
smooth functions and then by density,
\[
 \norm{v(b)}_{H^s}^2-\norm{v(a)}_{H^s}^2
 =2\int_a^b
 \pair{v_t(t)}{v(t)}_{H^{s-1},H^{s+1};H^s}\,dt.
\]
The shifted dual pairing obeys
\[
 \left|\pair{f}{g}_{H^{s-1},H^{s+1};H^s}\right|
 \leq\norm{f}_{H^{s-1}}\norm{g}_{H^{s+1}},
\]
so Cauchy--Schwarz proves \eqref{eq:LM-estimate}. The same approximation
argument gives weak \(H^s\)-continuity; weak continuity together with the
continuity of the squared norm supplied by the identity gives strong
\(H^s\)-continuity. Uniqueness follows because two continuous representatives
agreeing almost everywhere agree everywhere.
\end{proof}

Applying Lemma~\ref{lem:LM} to every temporal row gives, for all
\(N,m\in\N_0\),
\begin{equation}
 \partial_t^N u\in C([0,T_*];H^m).
 \label{eq:all-traces}
\end{equation}
Indeed, the two hypotheses are precisely adjacent coordinates of
\(\cI[u_0]\). This is the Lions--Magenes trace mechanism behind the endpoint,
whereas Lemma~\ref{lem:S0-endpoint} shows that the projected row
\(\Szero(I_*)\) already reaches the same endpoint through the equation.

\subsection{The adjacent critical block}

Put \(\Lambda=(-\Delta)^{1/2}\) and
\begin{equation}
 E_s(t):=\frac12\norm{\Lambda^s u(t)}_2^2,
 \qquad H(t):=E_{1/2}(t).
 \label{eq:energies}
\end{equation}

\begin{proposition}[Critical trace and current]
\label{prop:critical}
The adjacent coordinates in \(\cI[u_0]\) imply
\begin{align}
 u\in L^2(I_*;H^{5/2}),\quad
 u_t\in L^2(I_*;H^{1/2})
 &\Longrightarrow
 E_{3/2}\in L^\infty(I_*)\cap L^1(I_*),
 \label{eq:critical-first}\\
 u,u_t\in L^2(I_*;H^1)
 &\Longrightarrow H\in W^{1,1}(I_*).
 \label{eq:critical-second}
\end{align}
Moreover the pressure-explicit critical current
\begin{equation}
 \mathcal P_{1/2}
 :=-\pair{\Lambda^{1/2}u}
 {\Lambda^{1/2}\bigl((u\cdot\nabla)u+\nabla p\bigr)}
 \label{eq:critical-current}
\end{equation}
satisfies
\begin{equation}
 \mathcal P_{1/2}=H'+2\nu E_{3/2}\in L^1(I_*).
 \label{eq:critical-balance}
\end{equation}
\end{proposition}

\begin{proof}
Lemma~\ref{lem:LM} with \(s=3/2\) gives
\(u\in C([0,T_*];H^{3/2})\). More precisely, the same
time-mollification argument gives, for \(0\leq r\leq t\leq T_*\),
\begin{equation}
 E_{3/2}(t)-E_{3/2}(r)
 =\int_r^t\operatorname{Re}
 \bigl(\Lambda^{1/2}u_t(s),\Lambda^{5/2}u(s)\bigr)_2\,ds.
 \label{eq:critical-trace-identity}
\end{equation}
Cauchy--Schwarz therefore gives \(E_{3/2}\in W^{1,1}(I_*)\), and in
particular \(E_{3/2}\in L^\infty(I_*)\). Also
\[
 \int_0^{T_*}E_{3/2}(t)\,dt
 \leq\frac12\norm{u}_{L^2(I_*;H^{5/2})}^2<\infty,
\]
which proves \eqref{eq:critical-first}.

The second pair means \(u\in H^1(I_*;H^1)\). The Hilbert-space chain rule
therefore yields
\begin{equation}
 H'(t)=\pair{\Lambda^{1/2}u(t)}{\Lambda^{1/2}u_t(t)}
 \quad\text{for a.e. }t,
 \label{eq:H-prime}
\end{equation}
and
\[
 \norm{H'}_{L^1(I_*)}
 \leq\norm{u}_{L^2(I_*;H^1)}
       \norm{u_t}_{L^2(I_*;H^1)}<\infty.
\]
Thus \(H\in W^{1,1}(I_*)\).

On every compact subinterval of \(I_*\), pair \eqref{eq:NS} with
\(\Lambda u\). The viscous term is
\(-\nu\norm{\Lambda^{3/2}u}_2^2=-2\nu E_{3/2}\), so
\[
 H'=-2\nu E_{3/2}+\mathcal P_{1/2}.
\]
The pressure in \eqref{eq:critical-current} is the normalized solution of the
Poisson equation in \eqref{eq:pressure-projection}; equivalently, its global
pairing vanishes because \(\nabla\cdot\Lambda u=0\). The pressure-explicit
form is retained because it is the form differentiated in the jet below.
Since \(H'\) and \(E_{3/2}\) are both in \(L^1(I_*)\), the classical current
extends to the whole terminal interval as the \(L^1\) function in
\eqref{eq:critical-balance}.
\end{proof}

No additional critical-current estimate is required after the compatible
intersection is present: Proposition~\ref{prop:critical} is a trace consequence
of two adjacent rectangles.

\subsection{The unique projective endpoint}

For each \(m\), let \(u_*^{(m)}\in H^m\) denote the endpoint value furnished
by Lemma~\ref{lem:S0-endpoint}, equivalently by Lemma~\ref{lem:LM}. If
\(k>m\), then the \(H^k\)-continuous and \(H^m\)-continuous representatives
are the same function almost everywhere on \(I_*\); after embedding
\(H^k\hookrightarrow H^m\), continuity makes them identical on
\([0,T_*]\). In particular their endpoint values agree in \(H^m\).
Therefore the family \(\{u_*^{(m)}\}_m\) is compatible and defines one element
\begin{equation}
 u_*:=u(T_*)\in\bigcap_{m=0}^{\infty}H^m(\R^3)=\Hs(\R^3).
 \label{eq:endpoint}
\end{equation}
No bound uniform in \(m\) is asserted or needed. Since divergence is
continuous \(H^m\to H^{m-1}\), \(u_*\) is divergence-free; since the
\(H^0\)-trace is strong, it also has finite kinetic energy.

The same argument applied to \(\partial_t^N u\) in \eqref{eq:all-traces}
gives compatible endpoint traces
\begin{equation}
 V_N^*:=\lim_{t\uparrow T_*}\partial_t^Nu(t)\in\Hs,
 \qquad N\in\N_0.
 \label{eq:velocity-jet-traces}
\end{equation}

\subsection{The pressure-explicit endpoint jet}

For a multi-index \(\alpha\), set
\[
 J_{N,\alpha}:=\partial_t^N\partial_x^\alpha u,
 \qquad
 \Pi_{N,\alpha}:=\partial_t^N\partial_x^\alpha p.
\]
Direct differentiation of \eqref{eq:NS} gives, on \(I_*\),
\begin{equation}
\begin{aligned}
 J_{N+1,\alpha}
 &+
 \sum_{a=0}^{N}\sum_{\beta\leq\alpha}
 \binom Na\binom\alpha\beta
 (J_{a,\beta}\cdot\nabla)
 J_{N-a,\alpha-\beta}
 +\nabla\Pi_{N,\alpha}
 \\
 &=\nu\Delta J_{N,\alpha},
 \qquad \nabla\cdot J_{N,\alpha}=0,
\end{aligned}
\label{eq:full-binomial}
\end{equation}
while incompressibility fixes every pressure row through
\begin{equation}
 -\Delta\Pi_{N,\alpha}
 =\partial_i\partial_j\partial_t^N\partial_x^\alpha(u_i u_j).
 \label{eq:pressure-poisson-jet}
\end{equation}

For the pure-time jet write \(V_N=\partial_t^Nu\) and
\(Q_N=\partial_t^Np\). Equations
\eqref{eq:full-binomial}--\eqref{eq:pressure-poisson-jet} become
\begin{align}
 Q_N
 &=R_iR_j\sum_{a=0}^{N}\binom Na V_{a,i}V_{N-a,j},
 \label{eq:pressure-recurrence}\\
 V_{N+1}
 &=\nu\Delta V_N
 -\sum_{a=0}^{N}\binom Na(V_a\cdot\nabla)V_{N-a}
 -\nabla Q_N.
 \label{eq:velocity-recurrence}
\end{align}

\begin{proposition}[Endpoint jet reconstruction]
\label{prop:endpoint-jet}
The endpoint velocity \(u_*\) uniquely determines a sequence
\((V_N^*,Q_N^*)_{N\geq0}\subset\Hs\times\Hs\) by
\begin{align}
 V_0^*&=u_*,
 \label{eq:jet-start}\\
 Q_N^*&=R_iR_j\sum_{a=0}^{N}\binom Na V_{a,i}^*V_{N-a,j}^*,
 \label{eq:jet-pressure}\\
 V_{N+1}^*&=\nu\Delta V_N^*
 -\sum_{a=0}^{N}\binom Na(V_a^*\cdot\nabla)V_{N-a}^*
 -\nabla Q_N^*.
 \label{eq:jet-velocity}
\end{align}
This recursively constructed sequence equals the left endpoint traces of the
actual velocity-pressure jet.
\end{proposition}

\begin{proof}
The Fr\'echet space \(\Hs\) is closed under spatial differentiation and
pointwise multiplication. Riesz transforms preserve each \(H^m\), hence
preserve \(\Hs\). Thus \eqref{eq:jet-pressure}--\eqref{eq:jet-velocity}
construct one and only one \(\Hs\) pair at every order.

It remains to identify this formal recursion with the history. Fix \(m\).
The traces \eqref{eq:velocity-jet-traces} converge in arbitrarily high
Sobolev spaces. Sobolev multiplication, spatial differentiation, and the
Riesz transforms are continuous between the corresponding spaces. We may
therefore pass \(t\uparrow T_*\) in
\eqref{eq:pressure-recurrence}--\eqref{eq:velocity-recurrence}. Beginning
with \(V_0^*=u_*\), induction in \(N\) gives exactly
\eqref{eq:jet-pressure}--\eqref{eq:jet-velocity}. Since \(m\) was arbitrary,
the identities hold in \(\Hs\). Spatial differentiation then identifies the
full multi-index jet in \eqref{eq:full-binomial}.
\end{proof}

\subsection{Classical restart and contradiction of finite maximality}

We record the standard local theorem in precisely the form used here
\cite{Temam}.

\begin{lemma}[Smooth local restart and uniqueness]
\label{lem:local}
Let \(t_0\in\R\) and let \(a\in\Hs_\sigma(\R^3)\).  There is a
\(\delta>0\) and a unique smooth Navier--Stokes solution \((w,q)\) on
\([t_0,t_0+\delta]\times\R^3\) with \(w(t_0)=a\).  Every spatial Sobolev
order persists on this common interval, and the temporal jet at \(t_0\) is
the pressure-normalized recursion
\eqref{eq:jet-pressure}--\eqref{eq:jet-velocity} with \(V_0^*=a\).
\end{lemma}

\begin{proof}
This is the classical smooth-data local theorem for three-dimensional
Navier--Stokes \cite{Temam}.  Applied to \(a\in H^\infty\), it gives one
positive lifespan on which every higher Sobolev norm propagates along the same
solution.  Smoothness up to the initial face follows from the equation.
Differentiating there and using the normalized pressure Poisson equation gives
the stated jet recursion.  Uniqueness is the standard strong-solution
uniqueness on the same interval.
\end{proof}

Apply Lemma~\ref{lem:local} at time \(T_*\) with datum \(a=u_*\). It gives a
smooth solution \((w,q)\) on \([T_*,T_*+\delta]\). Define
\[
 (\widetilde u,\widetilde p)(t)
 :=
 \begin{cases}
 (u,p)(t),&0\leq t<T_*,\\
 (w,q)(t),&T_*\leq t\leq T_*+\delta.
 \end{cases}
\]
The velocity traces agree at \(T_*\) by construction. Therefore the
piecewise function has no distributional time-boundary defect and solves
\eqref{eq:NS} across \(T_*\). Proposition~\ref{prop:endpoint-jet} and the
last assertion of Lemma~\ref{lem:local} show that every left and right
velocity-pressure derivative agrees at \(T_*\). Hence
\((\widetilde u,\widetilde p)\) is classical, indeed smooth, on an interval
strictly larger than \([0,T_*)\).

Equivalently, one may start the local solution at a time
\(t_0<T_*\) sufficiently close to \(T_*\).  Continuity of the endpoint in
the fixed smooth-data local topology gives a common positive local lifespan,
and uniqueness identifies that solution with \(u\) on the overlap.  Both
descriptions restart the same datum-generated history.

This contradicts the definition of \(T_*\) as the maximal classical
lifespan. Thus the assumption \(T_*<\infty\) is impossible, and
\[
 \boxed{T_*=\infty.}
\]
Because the argument applies at every finite time and all velocity and
pressure rows lie in \(\Hs\), Sobolev embedding yields
\[
 \boxed{u,p\in C^\infty([0,\infty)\times\R^3).}
\]
This completes every implication in \eqref{eq:chain} and proves
Theorem~\ref{thm:main}.


\begin{thebibliography}{9}
\bibitem{LM}
J.-L. Lions and E. Magenes,
\emph{Non-Homogeneous Boundary Value Problems and Applications, Vol. I},
Springer, 1972.

\bibitem{Temam}
R. Temam,
\emph{Navier--Stokes Equations: Theory and Numerical Analysis},
AMS Chelsea, 2001.
\end{thebibliography}

\end{document}
